\section{Background and Related Work}
\label{sec:background}

Generative retrieval decodes an identifier for the target document instead of searching a dense index~\cite{tay2022dsi,decao2021genre}. A parallel line builds this identifier from a residual quantizer~\cite{lee2022rqvae}, learning \emph{discrete semantic IDs} instead of hand-built ones -- e.g.\ TIGER~\cite{rajput2023tiger} for recommendation -- though effectiveness degrades sharply with corpus size and task diversity~\cite{pradeep2023scale}. GENIUS~\cite{kim2025genius} is the multimodal instance: a CLIP-style encoder~\cite{radford2021clip} produces a joint image/text embedding, an RQ quantizes it into a semantic ID, and a T5~\cite{raffel2020t5} decoder generates a target ID via constrained beam search~\cite{decao2021genre} over M-BEIR~\cite{wei2024uniir}.

RQ codebooks are EMA-fit with no term rewarding even code usage, and are known to suffer \emph{codebook collapse}: a small code subset absorbs most assignments while the rest receive vanishing gradient~\cite{vandenoord2017vqvae}. The standard remedy adds an entropy-style auxiliary loss alongside reconstruction (\emph{balance regularization})~\cite{zhang2023regularizedvq}; an alternative line avoids collapse by construction instead, e.g.\ FSQ~\cite{mentzer2024fsq}. LETTER~\cite{wang2024letter} adds a diversity loss to RQ-VAE codebooks for generative recommendation -- the closest prior use of the intervention we study -- as one term in a joint objective for a single-modality identifier, not evaluated for its effect on downstream retrieval or on a codebook serving two target modalities at once, as ours does. None of this prior work evaluates whether balance regularization helps or hurts \emph{retrieval}, or whether its effect is stable across regimes -- the gap this paper addresses.
